\section{Discussion}

The preceding sections state what each safeguard does and what it refuses. This
section addresses a different question, asked repeatedly of the framing rather
than the mathematics: whether the presentation undersells the engine by
describing it mainly through what it prevents. The honest answer is
calibrated, not uniform: in one place the hedged framing is exactly right and
should not be loosened; in another, a genuinely nontrivial result is available
and is stated here explicitly.

\paragraph{What is a result, not only a guardrail.} The certified-multiplicity
boundary of Section~2.3(b) is not merely a safety mechanism sitting in front
of the classifier; it is itself a nontrivial, adversarially derived result.
Order-2 multiplicity is provably decidable from finite-precision evidence
alone, because two distinct eigenvalues never produce a nonzero rank-one
nullity at their mean; no analogous proof exists for order \(m>2\), and two
successive releases of this engine wrongly believed one did, each caught only
by an external counterexample (a linear-roundoff-coincidence carve-out that a
triangular matrix with genuinely distinct, non-roundoff-separated eigenvalues
defeated). The corrected boundary, that no floating-point evidence ever
certifies an order-\(m>2\) defect regardless of how small its members' spread
is, is verified against zero false certifications across perturbation scales
from \(10^{-2}\) to \(10^{-12}\) and orders 2 through 6, both directly and
under random orthogonal similarity. That is a stated, checkable claim about
what finite-precision multiplicity testing can and cannot establish in
general, not a disclaimer attached to a feature.

\paragraph{What remains a bound, not a theorem.} The constructive
counterexample of Section~3 (identical local reactant wells, activated rates
differing by \(e^{5}\approx148\) through the barrier alone) is exactly that: a
constructive counterexample under one stated model, sufficient to keep the
implementation from inferring rate from local architecture, and reproduced as
an executable regression test for that reason. It is not offered, and should
not be read, as a general limitation theorem for local-architecture-based rate
estimators; no such generality is proven. The distinction matters: a
counterexample bounds one specific inference for one stated construction, a
theorem would bound every construction of a stated class. Overstating the
former as the latter is the same error this framework spends Section~2.3 and
Section~4 refusing to make about degeneracy and about a simulated versus
measured anchor, and it is refused here for the same reason.

\paragraph{Reading the linewidth procedure and the worked examples.} The
three-step procedure of Part~III restates a standard decomposition; what it
adds is that each step is executable rather than declarative, refusing with a
stated reason when a step's precondition is not met (Section~5). The Rh(CO)\(_2\)acac
trajectory and the admission table of Section~6 demonstrate that procedure
running end to end on real systems, distinguishing what a source anchors from
what is illustrative at every step; the quantitative claim stays scoped to
the anchored endpoints, as the limitations that follow state plainly. Reframing either as a
general-purpose reproducible pipeline, independent of the specific systems
carried through it, would claim more than either demonstration establishes.
